\section{Choosing a smaller base}\label{g:reductions}
We now prove allowed-pole nonvanishing using the two metrics of
Section~\ref{g:analytic} and the weighted cotangent estimate of
Proposition~\ref{g:bt:weighted-slope}. The first two proofs apply when
there are no positive-degree holomorphic forms; the third gives a
different choice of base in the rationally connected case. All three
start from effective divisors with a fixed cotangent twist and must
produce both a horizontal section and every required adjoint rank.

\subsection{The common section conclusion}\label{rb:reductions}

\begin{theorem}[Allowed-pole nonvanishing]\label{g:allowed}\label{rb:allowed}
Let $V$ be a smooth connected projective variety over $\mathbb C$ with
$H^0(V,\Omega_V^j)=0$ for every $j>0$.  Let $D\geq0$ be integral and
suppose $-K_V+D$ carries a semipositive singular metric whose measure
$|s_D|^2$ has finite mass.  If a line bundle $L$ carries a semipositive
metric with bounded weights, then
\[
                    H^0(V,aL+cD)\ne0
\]
for some integers $a>0$ and $c\geq0$.  The same statement holds for a
rational line bundle $L$, with $a$ clearing its denominator.
\end{theorem}

In the three proofs below, whenever a dominant rational map from $V$
with connected general fibers is selected, use the model diagram of
Section~\ref{g:models}, with smooth projective base $Y$.
Thus $\bar p:\bar V\to V$ is birational,
$\bar f:\bar V\to Y$ is surjective and equidimensional with $\bar V$
normal, and a smooth resolution $W\to\bar V$ has composite maps
$p:W\to V$ and $f:W\to Y$. We write $J=K_{W/V}$.

\subsection{Maximal systems dominated by the target and boundary}
\label{rb:dominated}

\begin{proof}[First proof of Theorem~\ref{rb:allowed}]
Clear the denominator of $L$.  We induct on $\dim V$; a point is
immediate.  Call a nonempty linear system in a line bundle $P$
admissible if
\[
               aL+cD-P\quad\hbox{is pseudoeffective}
\]
for some integers $a>0,c\geq0$.  The constant system is admissible,
and products of admissible systems are admissible.  Choose one whose
image dimension is maximal.  If it is generically finite, then $P$ is
big.  Consequently $aL+cD$ is big, which supplies a section after a
positive multiple.  Otherwise, resolve its map and take the relative
algebraic closure of its image field to obtain $f:W\to Y$ with
connected general fibres and $\dim Y<\dim V$.

Every ratio of sections in every admissible system belongs to
$\mathbb C(Y)$.  Indeed adjoining such ratios cannot increase image
dimension, since products of the systems remain admissible.
The ratios are therefore algebraic over the chosen image field and
belong to its relative algebraic closure.  This argument also covers
subsystems and pencils.

Lemma~\ref{rb:determinant} gives $N_i\sim M+m_iL$, while
Proposition~\ref{rb:cotangent} gives $-M+kD$ pseudoeffective for some
$k\geq0$.  Systems in all positive combinations of the $N_i$ are
therefore admissible.  For $i<j<l$ the two effective divisors
\[
 (m_l-m_j)N_i+(m_j-m_i)N_l,
 \qquad (m_l-m_i)N_j
\]
are linearly equivalent, so the rational function relating them lies
in $\mathbb C(Y)$.  After pullback, their horizontal parts are equal.
It follows that the horizontal multiplicity at every prime is affine
in $m_i$.  Fixing $i<j$ and letting $l$ tend to infinity shows that its
slope is nonnegative. Thus $N_j-N_i$ defines a rational section $\rho$ of
$(m_j-m_i)L$ with no horizontal poles on $W$, and hence on $\bar V$.
If $Y$ is a point, it has no poles and already proves the theorem.

For a positive-dimensional base we verify the rank-one condition.
Let $P$ denote the original maximal system and let $A_0$ be the
pullback to $Y$ of its image hyperplane class.  It is big and
$p^*P-f^*A_0$ has an effective representative.  Given a sufficiently
ample base divisor $H$, a positive multiple of $A_0$ dominates $H$
up to an effective divisor.  Therefore $p_*f^*H$ is dominated in the
pseudoeffective order by a positive multiple of $L$ plus a nonnegative
multiple of $D$.

If $J+p^*D+mp^*L$ has two independent sections on a generic fibre,
spread them out and clear their vertical poles using $f^*H$ with $H$
sufficiently ample.  Their pushdowns are sections of
$D+mL+p_*f^*H$; the exceptional Jacobian $J$ contributes zero under
pushdown.  This system is admissible, but its ratio is nonconstant on
the generic fibre, contrary to the maximality property.  Hence its
generic dimension is at most one for every $m\geq0$.  For the multiples
of the exponent just constructed, $s_{J+p^*D}\rho^l$ supplies a nonzero
generic section, including $l=0$.  Lemmas~\ref{rb:rank-one} and
\ref{rb:transfer} apply.

The base has no positive-degree holomorphic forms: pullback by the
dominant map is injective, and forms on the smooth birational model
identify with forms on $V$.  The descended data thus satisfy exactly
the hypotheses of the theorem in lower dimension.  Induction gives a
section with poles on $D_Y$, and Lemma~\ref{rb:transfer} carries it to
a section with poles only on $D$.  An integral multiple of $D$ clears
those finitely many poles.
\end{proof}

\subsection{A supporting class and maximal combinations}
\label{rb:zero-degree}

Here is a different choice of the dimension-reducing map.  It uses
only effective combinations of the determinant divisors and the
permitted boundary, instead of all pseudoeffectively dominated
systems.

\begin{proof}[Second proof of Theorem~\ref{rb:allowed}]
Again induct on dimension and clear denominators.  If $L+D$ is big,
a positive multiple has a section.  Otherwise $L+D$ lies on the
boundary of the pseudoeffective cone.  Cone duality gives a nonzero
movable class $\alpha$ such that
\[
                       L\cdot\alpha=D\cdot\alpha=0.
\]
The zero-degree part of Proposition~\ref{rb:cotangent} and the determinant
construction give $N_i\sim M+m_iL$ with $M\cdot\alpha\leq0$.
Effectivity forces $N_i\cdot\alpha=M\cdot\alpha=0$.
No nonnegative integral combination of $D,N_1,N_2,\ldots$ is big.

Choose such a combination $G$ of maximal Iitaka dimension $d<\dim V$.
Take a multiple realizing dimension $d$, resolve its map, and take the
connected-fibre base $Y$.  Every ratio from the linear system of any
allowed combination lies in $\mathbb C(Y)$, since adding systems and
multiplying sections cannot increase $d$.  If $H$ is any base divisor,
then
\begin{equation}\label{rb:ample-domination}
                      ap^*G-f^*H\sim E_H\geq0
\end{equation}
for some sufficiently large $a$: the resolved system contains the
pullback of the image hyperplane class, whose transform on $Y$ is big.

We need the following stronger form of the ratio assertion.  A rational
function whose horizontal poles on $W$ are bounded by $p^*G'+E$, where
$G'$ is an allowed combination and $E$ is effective and $p$-exceptional,
belongs to $\mathbb C(Y)$.  Its remaining poles are vertical.  Choose
an effective sufficiently ample $H$ whose pullback bounds them; a
base divisor can contain the images of all these finitely many poles,
including images of codimension at least two, to the necessary orders.
By \eqref{rb:ample-domination}, the function is a ratio of sections of
$p^*(G'+aG)+E$.  Normality gives
$p_*\mathcal O_W(E)=\mathcal O_V$, so the exceptional term adds no
sections.  The original ratio assertion proves the claim.

The three-index relations among the $N_i$ now imply affine horizontal
multiplicities exactly as in the first proof.  Let $\rho$ have divisor
$N_j-N_i$ and exponent $k=m_j-m_i>0$.  It has no horizontal poles,
and its horizontal zeros are bounded by $p^*N_j$.  On a generic fibre
$s_{J+p^*D}\rho^l$ is a nonzero section of
$J+p^*D+lkp^*L$.  The ratio of any other section to this one has
horizontal poles bounded by
\[
                        J+p^*D+l p^*N_j.
\]
The stronger ratio assertion makes that ratio a base function.  Thus
the generic section space is one-dimensional for every $l\geq0$.
If $d=0$, $\rho$ is already regular.  Otherwise
Lemmas~\ref{rb:rank-one} and \ref{rb:transfer} give the two required
metrics and the controlled boundary on the smaller base.  Absence of
holomorphic forms descends, so dimension induction completes the
proof.
\end{proof}

\subsection{Maximal fibre spaces on a rationally connected variety}
\label{rb:max-fibration}

For rationally connected varieties one can choose a maximal fibre
space by two numerical and generic-section properties.  We give this
construction because its rank-one conclusion follows by enlarging a
field, rather than by fixing one class of systems in advance.

\begin{proposition}\label{rb:rc-method}
The conclusion of Theorem~\ref{rb:allowed} on a rationally connected
$V$ follows by dimension induction through rationally connected
bases, using fibre spaces that satisfy the following two conditions:
for an ample divisor $H$ on the base,
\[
 aL+cD-p_*f^*H\ \hbox{is pseudoeffective for some }a>0,c\geq0,
\]
and a positive multiple of $p^*L$ has a nonzero section on the generic
fibre without additional $D$-poles.
\end{proposition}

\begin{proof}
Suppose no positive multiple of $L$ has a section with poles only on
$D$.  In particular no $aL+cD$ with $a>0,c\geq0$ is big.
Rational connectedness gives absence of holomorphic forms, so take
the determinant divisors $N_i\sim M+m_iL$ and the inequality
$-M+kD$ pseudoeffective.  Let $K$ be the relative algebraic closure in
$\mathbb C(V)$ of the field generated by functions with divisors
\[
 (m_l-m_j)N_i+(m_j-m_i)N_l-(m_l-m_i)N_j,
                           \qquad i<j<l.
\]
These functions exist because the displayed divisor class is zero.
A finite selection gives the same transcendence degree, and the
relative algebraic closure is a finite extension of a finitely
generated subfield.  Thus $K$ defines a smooth projective base after
resolution.

Each of those functions is represented by a pencil whose total
bundle has class $uL+vM$, with $u,v>0$.  Finitely many such pencils give
a generically finite map from the base to their combined image; the
sum of the pulled-back hyperplane classes is therefore big on the
base.  The resolved pencils are bounded above by their total bundles.
Pushing down and using $-M+kD$ pseudoeffective proves the first
condition of the proposition.  A large multiple of the big base
class dominates any chosen ample divisor, so the same condition holds
with an ample divisor.

On the generic fibre all the displayed principal divisors vanish.
The effective restrictions of the $N_i$ therefore have affine
multiplicities in $m_i$.  Their slope is nonnegative by the unbounded
sequence, so $N_j-N_i$ supplies the second condition.  A point base
would already give a regular section, contrary to our supposition.
A base of full dimension would make the pushdown of an ample base
class big, hence would make some $aL+cD$ big, also impossible.
We have found a base of dimension strictly between zero and $\dim V$.

Choose, among fibre spaces with these two properties, one with maximal
base dimension.  The properties survive higher models.  For the first,
pullback of an ample class to a birational base model is big and a
large multiple dominates every fixed divisor; pulling and pushing
preserve effectivity.  The second is the regularity of a fixed rational
section over the generic fibre, which survives resolution.

Let $\rho$ be the supplied generic section of $kp^*L$.  If
$J+p^*D+mkp^*L$ had two independent generic sections, their ratio would
be transcendental over the relatively algebraically closed base field.
Adjoining it and taking relative algebraic closure enlarges the base
dimension.  Clear the two sections' vertical poles by a base twist
$f^*H'$.  Their ratio pencil has class
\[
                    J+p^*D+mkp^*L+f^*H'.
\]
Add an ample system from the old base.  The product map supplies a big
class on the enlarged base.  On pushdown $J$ vanishes, and the first
property for the old map bounds the remaining class by a multiple of
$L$ plus a multiple of $D$.  Thus the first property holds for the
enlarged map.  Restricting $\rho$ to its new generic fibre gives the
second property and stays nonzero at the generic point.  This
contradicts maximality.

The spaces in Lemma~\ref{rb:rank-one} consequently have dimension at
most one and have a nonzero element $s_{J+p^*D}\rho^m$, also at $m=0$.
Transfer gives the required data on the smaller base, which is
rationally connected because it is a smooth projective model of a
rational image of $V$.  Induction gives its permitted section and
transfer gives one on $V$, the required contradiction.
\end{proof}
